% 1 段組の論文（プレプリント体裁）。2 段組で起きる崩れの対照として要る。
% 段組を見ない実装に戻ったときでも 1 段組は通ってしまうので、ここが 100% でも
% 何の保証にもならない — そのことを可視化しておくための対照群。
%
% 正解の取り出し方は twocolumn.tex と同じ（本文 1 文 = \s{...} を 1 行）。
\documentclass[11pt,a4paper]{article}

\usepackage{mathptmx}
\usepackage[margin=2.4cm]{geometry}
\usepackage{amsmath}
\usepackage{fancyhdr}
\usepackage{booktabs}

\newcommand{\s}[1]{#1}

\pagestyle{fancy}
\fancyhf{}
\fancyhead[L]{Sparse Attention Notes}
\fancyfoot[C]{\thepage}

\title{Notes on Selecting Attention Blocks}
\author{Mei Tanaka}
\date{}

\begin{document}
\maketitle

\begin{abstract}
\s{These notes collect the design decisions behind the block selection rule.}
\s{They are written for readers who have already seen the main paper and want the details.}
\s{We cover the choice of block size, the treatment of padding and the cost of the selection pass.}
\end{abstract}

\section{Block Size}

\s{The block size controls the trade off between the cost of selection and the granularity of the pattern.}
\s{Small blocks approximate dense attention closely but make the selection matrix large.}
\s{Large blocks make the selection cheap and force the kernel to read positions it does not need.}
\s{We measured both effects on a single node with eight accelerators.}
\s{A block size of sixty four was the best compromise on that hardware.}

\s{The optimal block size depends on the memory system rather than on the model.}
\s{On a machine with a larger cache the optimum moved to one hundred and twenty eight.}
\s{We therefore expose the block size as a configuration value rather than hard coding it.}

\section{Padding}

\s{Sequences rarely divide evenly into blocks, so the last block is padded.}
\s{Padded positions are given a summary of zero and a score of negative infinity.}
\s{This keeps the padded block out of the retained set without a separate branch in the kernel.}
\s{The same trick is used for the causal mask inside a retained block.}

\s{The cost of the selection pass is the product of the number of query blocks and key blocks.}
\begin{equation}
  C = \frac{n^2}{b^2} \, d ,
  \label{eq:cost}
\end{equation}
\s{Here n is the sequence length, b is the block size and d is the head dimension.}

\begin{table}[t]
  \centering
  \begin{tabular}{lrr}
    \toprule
    Block size & Selection cost & Accuracy \\
    \midrule
    16 & 12.4 & 88.9 \\
    64 & 1.9 & 88.7 \\
    256 & 0.4 & 85.1 \\
    \bottomrule
  \end{tabular}
  \caption{Cost of the selection pass in percent of a dense attention call, and accuracy on the held out split.}
  \label{tab:block}
\end{table}

\section{Discussion}

\s{The selection rule is simple enough that it can be reimplemented in an afternoon.}
\s{Most of the engineering effort went into the fused kernel rather than into the rule itself.}
\s{We expect better summaries to matter more than a better selection rule \cite{tanaka2023blocks}.}
\s{A learned summary would add parameters, which we wanted to avoid for this release.}

\begin{thebibliography}{9}
\bibitem{tanaka2023blocks}
M. Tanaka. Block structured kernels. In Proc. ICML, pages 991--1002, 2023.
\bibitem{roe2025efficient}
R. Roe. Efficient attention in practice. Journal of ML, 12(3):45--67, 2025.
\end{thebibliography}

\appendix

\section{Measurement Protocol}

\s{Every number in these notes is the mean of five runs.}
\s{We discard the first run of each configuration because it includes kernel compilation.}
\s{The machine was otherwise idle and the clocks were locked to their base frequency.}

\end{document}
